\documentclass[12pt,a4paper]{article}

\usepackage[utf8]{inputenc}
\usepackage[russian]{babel}
\usepackage{amsmath, amsfonts, amssymb, amsthm}
\usepackage{geometry}
\geometry{top=2cm, bottom=2cm, left=2.5cm, right=1.5cm}
\usepackage{booktabs}
\usepackage{graphicx}
\usepackage{xcolor}

\theoremstyle{definition}
\newtheorem{definition}{Определение}
\newtheorem{example}{Пример}
\newtheorem{algorithm}{Алгоритм}

\theoremstyle{plain}
\newtheorem{theorem}{Теорема}

\title{Билет №74 \\ \small Нормальные формы схемы БД с первой по пятую, нормальная форма Бойса-Кодда. (СпБГУ)}
\author{Сериков Владислав Александрович}
\date{16 мая 2026}

\begin{document}

\maketitle
\tableofcontents
\newpage

\section{Введение и базовые определения}

Для строгого понимания нормальных форм необходимо ввести базовые понятия реляционной алгебры и теории зависимостей.

\begin{definition}[Отношение и схема]
\textbf{Схема отношения} $R(A_1, A_2, \dots, A_n)$ — это конечное множество атрибутов.
\textbf{Отношение} $r$ со схемой $R$ — это конечное множество кортежей, где каждый кортеж представляет собой отображение атрибутов в их домены.
\end{definition}

\begin{definition}[Функциональная зависимость (ФЗ)]
Пусть $X$ и $Y$ — подмножества атрибутов схемы $R$. Говорят, что $Y$ \textbf{функционально зависит} от $X$ (обозначается $X \to Y$), если для любых двух кортежей $t_1, t_2 \in r$ из того, что $t_1[X] = t_2[X]$, следует $t_1[Y] = t_2[Y]$.
\end{definition}

\begin{definition}[Ключи]
\textbf{Суперключ} — это подмножество атрибутов $K \subseteq R$, такое что $K \to R$. \\
\textbf{Потенциальный ключ (Candidate Key)} — это минимальный (по включению) суперключ. \\
Атрибут называется \textbf{первичным (prime)}, если он входит хотя бы в один потенциальный ключ отношения. В противном случае атрибут является \textbf{непервичным (non-prime)}.
\end{definition}

\section{Теорема Хита и декомпозиция без потерь}

Центральным понятием нормализации является замена одного отношения на несколько других без потери информации (свойство соединения без потерь).

\begin{theorem}[Теорема Хита / Heath's Theorem]
Пусть $R(X, Y, Z)$ — схема отношения, где $X, Y, Z$ — непересекающиеся множества атрибутов, и $X \cup Y \cup Z = R$. Если в отношении $r(R)$ выполняется функциональная зависимость $X \to Y$, то декомпозиция $R$ на $R_1(X, Y)$ и $R_2(X, Z)$ является декомпозицией без потерь, то есть:
$$r = \pi_{X, Y}(r) \bowtie \pi_{X, Z}(r)$$
где $\pi$ — операция проекции, а $\bowtie$ — операция естественного соединения.
\end{theorem}

\begin{proof}
Для доказательства равенства множеств необходимо показать взаимное включение: $r \subseteq \pi_{X, Y}(r) \bowtie \pi_{X, Z}(r)$ и $\pi_{X, Y}(r) \bowtie \pi_{X, Z}(r) \subseteq r$.

\textbf{Часть 1: $r \subseteq \pi_{X, Y}(r) \bowtie \pi_{X, Z}(r)$} \\
Это свойство выполняется для любой декомпозиции по определению операций проекции и соединения. Если кортеж $t \in r$, то $t[X, Y] \in \pi_{X, Y}(r)$ и $t[X, Z] \in \pi_{X, Z}(r)$. При естественном соединении по общему атрибуту $X$ эти проекции соединятся, породив исходный кортеж $t$.

\textbf{Часть 2: $\pi_{X, Y}(r) \bowtie \pi_{X, Z}(r) \subseteq r$} \\
1. Пусть кортеж $t \in \pi_{X, Y}(r) \bowtie \pi_{X, Z}(r)$. Нам нужно показать, что $t \in r$. \\
2. По определению естественного соединения, существуют кортежи $t_1 \in \pi_{X, Y}(r)$ и $t_2 \in \pi_{X, Z}(r)$, такие что $t_1[X] = t_2[X] = t[X]$, $t_1[Y] = t[Y]$ и $t_2[Z] = t[Z]$. \\
3. Так как $t_1 \in \pi_{X,Y}(r)$, в исходном отношении $r$ существует некоторый кортеж $u$, такой что $u[X] = t[X]$ и $u[Y] = t[Y]$. \\
4. Аналогично, из-за $t_2 \in \pi_{X,Z}(r)$, в $r$ существует кортеж $v$, такой что $v[X] = t[X]$ и $v[Z] = t[Z]$. \\
5. Мы имеем два кортежа $u, v \in r$, у которых совпадают значения атрибутов $X$: $u[X] = v[X]$. \\
6. По условию теоремы выполняется функциональная зависимость $X \to Y$. Следовательно, совпадение $X$ влечет за собой совпадение $Y$: $u[Y] = v[Y]$. \\
7. Из шага 3 мы знаем, что $u[Y] = t[Y]$. Значит, и $v[Y] = t[Y]$. \\
8. Рассмотрим кортеж $v$. Для него выполняется: $v[X] = t[X]$ (шаг 4), $v[Y] = t[Y]$ (шаг 7), $v[Z] = t[Z]$ (шаг 4). \\
9. Поскольку схема отношения состоит только из атрибутов $X, Y, Z$, кортеж $v$ полностью совпадает с кортежем $t$. \\
10. Так как $v \in r$, то и $t \in r$. Включение доказано. $\blacksquare$
\end{proof}


\section{Нормальные формы на основе функциональных зависимостей (1NF - 3NF, BCNF)}

\subsection{Первая нормальная форма (1NF)}
\begin{definition}
Отношение находится в \textbf{1NF}, если домен каждого атрибута содержит только атомарные (неделимые) значения, и каждый кортеж содержит ровно одно значение для каждого атрибута.
\end{definition}

\begin{example}
\textbf{Нарушение 1NF:} Отношение \texttt{Студент(ID, Имя, Телефоны)}, где в поле \texttt{Телефоны} записано «123-45, 678-90». \\
\textbf{Приведение к 1NF:} Создание двух отдельных строк для каждого телефона студента или вынос телефонов в отдельную таблицу \texttt{Телефоны(Студент\_ID, Телефон)}.
\end{example}

\subsection{Вторая нормальная форма (2NF)}
\begin{definition}
Отношение находится во \textbf{2NF}, если оно находится в 1NF и каждый непервичный атрибут \textit{полностью функционально зависит} от каждого потенциального ключа (т.е. не существует зависимости непервичного атрибута от части составного ключа).
\end{definition}

\begin{example}
Пусть $R(\underline{\text{Студент}}, \underline{\text{Курс}}, \text{Оценка}, \text{Преподаватель})$. Ключ: $\{\text{Студент, Курс}\}$. \\
ФЗ: $\{\text{Студент, Курс}\} \to \text{Оценка}$; $\text{Курс} \to \text{Преподаватель}$. \\
\textbf{Нарушение:} Атрибут «Преподаватель» зависит только от части ключа («Курс»). \\
\textbf{Приведение к 2NF:} $R_1(\underline{\text{Студент, Курс}}, \text{Оценка})$ и $R_2(\underline{\text{Курс}}, \text{Преподаватель})$.
\end{example}

\subsection{Третья нормальная форма (3NF)}
\begin{definition}
Отношение находится в \textbf{3NF}, если оно находится во 2NF и никакой непервичный атрибут не зависит транзитивно от потенциального ключа. Формально: для любой нетривиальной ФЗ $X \to A$ выполняется хотя бы одно из условий:
1) $X$ — суперключ; \\
2) $A$ — первичный атрибут.
\end{definition}

\begin{example}
$R(\underline{\text{Сотрудник}}, \text{Отдел}, \text{Локация\_Отдела})$. \\
ФЗ: $\text{Сотрудник} \to \text{Отдел}$, $\text{Отдел} \to \text{Локация\_Отдела}$. \\
\textbf{Нарушение:} Локация зависит от ключа транзитивно через Отдел. Зависимость $\text{Отдел} \to \text{Локация\_Отдела}$ нарушает 3NF, так как Отдел не суперключ, а Локация — непервичный атрибут. \\
\textbf{Приведение к 3NF:} $R_1(\underline{\text{Сотрудник}}, \text{Отдел})$ и $R_2(\underline{\text{Отдел}}, \text{Локация\_Отдела})$.
\end{example}

\subsection{Нормальная форма Бойса-Кодда (BCNF)}
\begin{definition}
Отношение находится в \textbf{BCNF}, если для любой нетривиальной функциональной зависимости $X \to A$ левая часть $X$ является суперключом.
\end{definition}
\textit{Замечание:} BCNF строже 3NF, так как не делает исключений для первичных атрибутов.

\begin{example}
$R(\text{Студент}, \text{Предмет}, \text{Преподаватель})$. Каждый преподаватель ведет только один предмет. Студент может слушать предмет только у одного преподавателя. \\
ФЗ: $\{\text{Студент, Предмет}\} \to \text{Преподаватель}$; $\text{Преподаватель} \to \text{Предмет}$. \\
Ключи: $\{\text{Студент, Предмет}\}$, $\{\text{Студент, Преподаватель}\}$. \\
Все атрибуты — первичные. Отношение находится в 3NF. \\
\textbf{Нарушение BCNF:} В зависимости $\text{Преподаватель} \to \text{Предмет}$ детерминант «Преподаватель» \textbf{не} является суперключом.
\end{example}

\vspace{1em}
\begin{algorithm}[Декомпозиция отношения в BCNF без потерь]
\textbf{Вход:} Схема отношения $R$, множество ФЗ $F$. \\
\textbf{Выход:} Множество схем $\{R_1, \dots, R_k\}$, находящихся в BCNF и сохраняющих информацию (по теореме Хита).

\textbf{Шаги алгоритма:}
\begin{enumerate}
    \item Инициализировать результирующее множество схем: $\mathcal{D} = \{R\}$.
    \item Пока в $\mathcal{D}$ существует схема $S$, которая не находится в BCNF по отношению к $F$:
    \begin{enumerate}
        \item Найти нетривиальную ФЗ $X \to Y$, которая выполняется на $S$ и нарушает BCNF (т.е. $X \cap Y = \emptyset$ и $X$ не является суперключом в $S$).
        \item Заменить схему $S$ в множестве $\mathcal{D}$ на две новые схемы: $S_1 = X \cup Y$ и $S_2 = S \setminus Y$.
    \end{enumerate}
    \item Возвратить $\mathcal{D}$.
\end{enumerate}
\end{algorithm}

\textbf{Иллюстрация шагов к Примеру (BCNF):}
\begin{itemize}
    \item Вход: $R(\text{Студент}, \text{Предмет}, \text{Преп})$, ФЗ: $\text{Преп} \to \text{Предмет}$.
    \item Шаг 1: $\mathcal{D} = \{R(\text{Студент}, \text{Предмет}, \text{Преп})\}$.
    \item Шаг 2: В $R$ ФЗ $\text{Преп} \to \text{Предмет}$ нарушает BCNF.
    \item Шаг 2(b): Разбиваем на $S_1 = \{\text{Преп, Предмет}\}$ и $S_2 = \{\text{Студент, Предмет, Преп}\} \setminus \{\text{Предмет}\} = \{\text{Студент, Преп}\}$.
    \item Шаг 3: $\mathcal{D} = \{S_1(\text{Преп, Предмет}), S_2(\text{Студент, Преп})\}$. Обе схемы в BCNF.
\end{itemize}

\section{Многозначные зависимости и 4NF}

\begin{definition}[Многозначная зависимость (МЗ)]
Говорят, что на отношении $R$ выполняется \textbf{многозначная зависимость} $X \twoheadrightarrow Y$, если для любых двух кортежей $t_1, t_2 \in r$, совпадающих по $X$ ($t_1[X] = t_2[X]$), в отношении $r$ обязательно существует кортеж $t_3$, такой что $t_3[X] = t_1[X]$, $t_3[Y] = t_1[Y]$ и $t_3[R \setminus (X \cup Y)] = t_2[R \setminus (X \cup Y)]$. \\
Интуитивно: множество значений $Y$, ассоциированных с данным значением $X$, не зависит от значений остальных атрибутов отношения.
\end{definition}

\begin{definition}[Четвертая нормальная форма (4NF)]
Отношение находится в \textbf{4NF}, если оно находится в BCNF и для любой нетривиальной многозначной зависимости $X \twoheadrightarrow Y$ множество $X$ является суперключом.
\end{definition}

\begin{example}
$R(\text{Ресторан}, \text{Блюдо}, \text{Район\_Доставки})$. \\
Ресторан готовит множество блюд и осуществляет доставку во множество районов. Меню блюд и районы доставки не зависят друг от друга. \\
МЗ: $\text{Ресторан} \twoheadrightarrow \text{Блюдо}$ и $\text{Ресторан} \twoheadrightarrow \text{Район\_Доставки}$. \\
\textbf{Нарушение 4NF:} Ключом является вся строка целиком $\{\text{Ресторан, Блюдо, Район\_Доставки}\}$. Детерминант МЗ «Ресторан» не суперключ. Возникает избыточность (декартово произведение блюд и районов для каждого ресторана). \\
\textbf{Приведение к 4NF:} Декомпозиция на $R_1(\text{Ресторан}, \text{Блюдо})$ и $R_2(\text{Ресторан}, \text{Район\_Доставки})$.
\end{example}

\section{Зависимости соединения и 5NF}

\begin{definition}[Зависимость соединения (ЗС)]
Отношение $R$ удовлетворяет \textbf{зависимости соединения} $\bowtie(R_1, R_2, \dots, R_n)$, где $R_1 \cup R_2 \cup \dots \cup R_n = R$, если отношение $R$ всегда равно естественному соединению своих проекций на $R_1, \dots, R_n$:
$$R = \pi_{R_1}(R) \bowtie \pi_{R_2}(R) \bowtie \dots \bowtie \pi_{R_n}(R)$$
\end{definition}

\begin{definition}[Пятая нормальная форма (5NF / PJNF)]
Отношение находится в \textbf{5NF} (или Project-Join Normal Form), если оно находится в 4NF и каждая нетривиальная зависимость соединения $\bowtie(R_1, \dots, R_n)$ вытекает из потенциальных ключей отношения (то есть каждый $R_i$ содержит суперключ исходного отношения).
\end{definition}

\begin{example}
Рассмотрим отношение $R(\text{Агент}, \text{Компания}, \text{Товар})$. \\
Смысл отношения: Агент представляет Компанию; Компания производит Товар; Агент продает Товар. \\
Существует правило: \textit{Если Агент представляет Компанию, и Компания производит Товар, и этот же Агент продает данный Товар, то соответствующая строка (Агент, Компания, Товар) должна быть в БД.} \\
Это круговое условие порождает зависимость соединения:
$$\bowtie(\{\text{Агент, Компания}\}, \{\text{Компания, Товар}\}, \{\text{Агент, Товар}\})$$
\textbf{Нарушение 5NF:} Отношение не может быть декомпозировано без потерь на две таблицы, но может быть декомпозировано на три: $R_1(\text{Агент}, \text{Компания})$, $R_2(\text{Компания}, \text{Товар})$, $R_3(\text{Агент}, \text{Товар})$. Ни одна из этих схем не содержит полный суперключ оригинала. \\
В 5NF исходное отношение заменяется тремя указанными бинарными отношениями, что исключает аномалии обновления, возникающие при добавлении или удалении связи между Агентом, Компанией и Товаром.
\end{example}

\end{document}
